% GENERATED FILE. Edit canonical lessons, then run npm run generate:books.
% canonical-source-hash: fnv1a64:458231ca7b8a0572
% canonical-lessons: MW-C162-jasu, MW-C162-koni-jasu, MW-C162-kal, MW-R162-tomorrow

\chapter{Tomorrow}
\label{ch:mw-a2-tomorrow}
\hlchaptermodality{162}

\begin{chapteropening}
\textbf{By the end of this chapter:} \emph{I can say what I will and will not do, and tell yesterday from tomorrow by the verb.}

Complete MW-R162-tomorrow: say what you will and won't do, and use \mw{काल} for yesterday and tomorrow.
\end{chapteropening}

\section[\texorpdfstring{\mw{जासूं} (jāsū̃)}{जासूं (jāsū̃)}]{\mw{हूं} \mw{काल} \mw{जासूं} — I will go}

\label{lesson:MW-C162-jasu}



\begin{quote}
Say \emph{I went to the market yesterday}. The same \textbf{\mw{काल}} also points the other way, to tomorrow.
\end{quote}

\begin{grammarlens}[title={Grammar Lens: -\mw{सूं}}]
Add \textbf{-\mw{सूं}} to a stem and you have \emph{I will}:

\begin{quote}
\textbf{\mw{हूं} \mw{काल} \mw{बाजार} \mw{जासूं।}}
\end{quote}

\begin{quote}
\emph{hū̃ kāl bājār jāsū̃.} — I will go to the market tomorrow.
\end{quote}

No \textbf{\mw{हूं}} at the end: the future stands on its own. Every stem takes it the same way: \textbf{\mw{आसूं}}, \textbf{\mw{खासूं}}, \textbf{\mw{बोलसूं}}, \textbf{\mw{समझसूं}}. CDIAL quotes an Old Marwari line with this future in it, \textquotedblleft{}I will fan\textquotedblright{} (CDIAL 5593). The spelling \textbf{-\mw{सूं}} is central Marwari as this book writes it; other districts differ.
\end{grammarlens}

\subsection*{Your turn}
You know two set phrases that look ahead: \textbf{\mw{पाछे} \mw{मिलसू}} and \textbf{\mw{कितणो} \mw{आखरी} \mw{भाव} \mw{दोगे}?} Keep saying them as you learned them.

Say these aloud:

\begin{itemize}
\raggedright
  \item \textbf{\mw{हूं} \mw{काल} \mw{आसूं।}} — \emph{hū̃ kāl āsū̃.} — I will come tomorrow.
  \item \textbf{\mw{हूं} \mw{आथण} \mw{दूध} \mw{पीसूं।}} — \emph{hū̃ āthaṇ dūdh pīsū̃.} — I will drink milk in the evening.
  \item \textbf{\mw{हूं} \mw{काल} \mw{बाजार} \mw{जासूं।}} — \emph{hū̃ kāl bājār jāsū̃.} — I will go to the market tomorrow.
\end{itemize}

\begin{tcolorbox}[breakable,colback=teal!4,colframe=teal!35!black,title={Before you move on}]
\emph{I will go}? (\textbf{\mw{जासूं}}.) \emph{I will come}? (\textbf{\mw{आसूं}}.) What comes after it? (Nothing.)
\end{tcolorbox}
\section[\texorpdfstring{\mw{कोनी} \mw{जासूं} (konī jāsū̃)}{कोनी जासूं (konī jāsū̃)}]{\mw{हूं} \mw{काल} \mw{कोनी} \mw{जासूं} — I won't go}

\label{lesson:MW-C162-koni-jasu}



\begin{quote}
Say \emph{I will come tomorrow}. Then \emph{I don't drink tea}.
\end{quote}

\begin{grammarlens}[title={Grammar Lens: \mw{कोनी} before the future}]
\textbf{\mw{कोनी}} goes in front of the future, as it does in front of the present:

\begin{quote}
\textbf{\mw{हूं} \mw{काल} \mw{कोनी} \mw{जासूं।}}
\end{quote}

\begin{quote}
\emph{hū̃ kāl konī jāsū̃.} — I won't go tomorrow.
\end{quote}

\begin{quote}
\textbf{\mw{हूं} \mw{परसूं} \mw{कोनी} \mw{आसूं।}}
\end{quote}

\begin{quote}
\emph{hū̃ parsū̃ konī āsū̃.} — I won't come the day after tomorrow.
\end{quote}
\end{grammarlens}

\subsection*{Your turn}
Say these aloud:

\begin{itemize}
\raggedright
  \item \textbf{\mw{हूं} \mw{मांस} \mw{कोनी} \mw{खासूं।}} — \emph{hū̃ mā̃s konī khāsū̃.} — I won't eat meat.
  \item \textbf{\mw{हूं} \mw{आज} \mw{बाजार} \mw{कोनी} \mw{जासूं।}} — \emph{hū̃ āj bājār konī jāsū̃.} — I won't go to the market today.
\end{itemize}

\begin{tcolorbox}[breakable,colback=teal!4,colframe=teal!35!black,title={Before you move on}]
\emph{I won't go}? (\textbf{\mw{कोनी} \mw{जासूं}}.) \emph{I don't go}? (\textbf{\mw{कोनी} \mw{जाऊं}}.)
\end{tcolorbox}
\section[\texorpdfstring{\mw{काल} (kāl)}{काल (kāl)}]{\mw{काल} — yesterday or tomorrow}

\label{lesson:MW-C162-kal}



\begin{quote}
Say \emph{I won't go tomorrow}. Then say \textbf{\mw{काल}} and \textbf{\mw{परसूं}} on their own.
\end{quote}

\begin{grammarlens}[title={Grammar Lens: the verb decides}]
\textbf{\mw{काल}} is one word for two days, and the verb tells you which:

\begin{quote}
\textbf{\mw{हूं} \mw{काल} \mw{गयो।}}
\end{quote}

\begin{quote}
\emph{hū̃ kāl gayo.} — I went yesterday.
\end{quote}

\begin{quote}
\textbf{\mw{हूं} \mw{काल} \mw{जासूं।}}
\end{quote}

\begin{quote}
\emph{hū̃ kāl jāsū̃.} — I will go tomorrow.
\end{quote}

\textbf{\mw{परसूं}} works the same way, for the day before yesterday or the day after tomorrow.
\end{grammarlens}

\subsection*{Your turn}
Say these aloud:

\begin{itemize}
\raggedright
  \item \textbf{\mw{राम} \mw{काल} \mw{आयो।}} — \emph{rām kāl āyo.} — Ram came yesterday.
  \item \textbf{\mw{हूं} \mw{काल} \mw{आसूं।}} — \emph{hū̃ kāl āsū̃.} — I will come tomorrow.
  \item \textbf{\mw{मां} \mw{परसूं} \mw{आई।}} — \emph{mā̃ parsū̃ āī.} — Mother came the day before yesterday.
\end{itemize}

\begin{tcolorbox}[breakable,colback=teal!4,colframe=teal!35!black,title={Before you move on}]
In \textbf{\mw{हूं} \mw{काल} \mw{गयो।}}, which day? (Yesterday.) In \textbf{\mw{हूं} \mw{काल} \mw{जासूं।}}? (Tomorrow.)
\end{tcolorbox}
\section[(dialogue)]{Tomorrow — a review}

\label{lesson:MW-R162-tomorrow}



\begin{quote}
Say \emph{I will come tomorrow}.
\end{quote}

\subsection*{Your turn}
Say these aloud:

\begin{itemize}
\raggedright
  \item \emph{I will go to the market tomorrow}; \emph{I won't go to the market today}
  \item \emph{I will drink milk in the evening}; \emph{I won't eat meat}
  \item \emph{I went yesterday}; \emph{I will go tomorrow}
  \item \emph{I drew water}; \emph{I will draw water tomorrow}
\end{itemize}

\begin{tcolorbox}[breakable,colback=teal!4,colframe=teal!35!black,title={Before you move on}]
\emph{I will go}? (\textbf{\mw{जासूं}}.) \emph{I won't go}? (\textbf{\mw{कोनी} \mw{जासूं}}.) \emph{I went}? (\textbf{\mw{गयो}}.)
\end{tcolorbox}
